% =====================================================================
%  CHAPTER 12: 学校突发公共卫生事件应急管理（对应大纲第十一章）
%  参考PPT：第十八章 学校突发公共卫生事件应急管理（李生慧，上海交通大学）
% =====================================================================
\chapter{学校突发公共卫生事件应急管理}
\setcounter{chapter}{12}
\chapterintro{
掌握学校突发公共卫生事件的定义、类型和分级；掌握学校传染病事件的定义、流行病学特点和应急处置（疫情报告、调查、控制、信息发布、应急反应终止）；掌握学校食源性疾病、食物中毒和食品安全事故的概念；掌握学校食物中毒事故的类型、流行特征、食品安全监督与处置；掌握群体心因性反应事件的定义、流行特征、早期识别和应急处置；熟悉学校突发公共卫生事件应对的工作原则和措施。
}

% ----- 12.1 概述 -----
\section{学校突发公共卫生事件概述}

\begin{defbox}
\textbf{学校突发公共卫生事件（Public Health Emergency in School）：}是指在学校内突然发生，造成或可能造成师生员工健康严重损害的重大传染病疫情、重大食物或职业中毒、不明原因的群体性疾病、群体性异常反应以及其他严重影响师生员工身心健康的公共卫生事件。
\end{defbox}

\subsection{学校突发公共卫生事件的类型}

\begin{keybox}
较为常见的类型主要有：
\begin{enumerate}[leftmargin=*]
    \item \imp{重大传染病疫情}
    \item \imp{预防接种和预防服药群体性不良事件}
    \item \imp{群体性不明原因疾病}
    \item \imp{食物中毒}
    \item \imp{其他中毒}
    \item \imp{环境因素事件}
    \item \imp{意外辐射照射事件}
\end{enumerate}
\end{keybox}

\subsection{学校突发公共卫生事件的分级}

\begin{defbox}
根据《国家突发公共卫生事件应急预案》（2006年）分为四级：
\begin{enumerate}[leftmargin=*]
    \item \imp{特别重大突发公共卫生事件（I级）}
    \item \imp{重大突发公共卫生事件（II级）}
    \item \imp{较大突发公共卫生事件（III级）}
    \item \imp{一般突发公共卫生事件（IV级）}
\end{enumerate}
\end{defbox}

\subsection{应对学校突发公共卫生事件的工作原则和措施}

\begin{keybox}
\textbf{（一）政府层面的应对：}建立应急组织和指挥机构；建立专家咨询委员会；建立应急处理专业技术机构。

\textbf{（二）教育行政部门的应对：}建立应急日常管理组织机构；制订应急相应预案；启动应急反应措施；进行事件信息发布。

\textbf{（三）学校层面的应对：}强化单位领导负责制；健全学校各项健康管理制度；严格落实事件报告管理制度；配合卫生部门采取有效措施；加强学生健康教育。
\end{keybox}

% ----- 12.2 传染病事件 -----
\section{学校传染病事件应对}

\subsection{学校传染病事件的类型和流行病学特点}

\begin{defbox}
\textbf{学校传染病事件（Infectious Disease Events in School）：}是指学校某种传染病在短时间内发生、波及范围广泛，出现大量的病人或死亡病例，发病率远远超过常年发病率水平的紧急情况。\imp{传染病事件是学校突发公共卫生事件最主要的类型，监测资料显示其所占比例高达90\%$\sim$95\%，且以呼吸道和消化道传染病暴发疫情为主。}

\textbf{学校传染病事件的类型：}(1)《中华人民共和国传染病防治法》规定的法定传染病；(2)国务院卫生行政部门决定并公布列入乙类、丙类传染病的其他传染病；(3)省、自治区、直辖市人民政府决定并公布的按照乙类、丙类传染病管理的其他传染病；(4)新发传染病（全球首次发现）；(5)我国尚未发现的传染病（如埃博拉、猴痘、黄热病等）；(6)我国已消灭的传染病（天花、脊髓灰质炎）。
\end{defbox}

\begin{keybox}
\textbf{学校传染病事件的流行病学特点：}

\begin{enumerate}[leftmargin=*]
    \item \imp{多以呼吸道和消化道传染病为主}：调查显示，2014$\sim$2016年学校报告的传染病突发公共卫生事件主要是呼吸道传染病，其次是肠道传染病，两者占总事件的97.91\%。
    \item \imp{存在地区和学校分布差异}：社会经济相对落后的西南省份较为多发；乡村中、小学和县城小学报告的传染病暴发事件最多。
    \item \imp{全年两个发病高峰}：发生时间呈现双峰分布，高峰分别在每年的3$\sim$6月和10$\sim$12月之间。
\end{enumerate}
\end{keybox}

\subsection{学校传染病的监督管理与应急处置}

\begin{keybox}
\textbf{（一）学校传染病监督管理}

依照相关法律法规，学校应从如下方面加强传染病的监督管理和能力建设：

\begin{enumerate}[leftmargin=*]
    \item \imp{建设卫生管理制度}：制定和完善学校各项卫生管理制度，配套相应的实施方案和监督机制，保证经费和人员的匹配，定期组织专业培训。
    \item \imp{标准化卫生工作流程}：细化各项卫生工作的规范标准和具体操作流程，包括加强饮用水和食堂卫生管理，完善洗手设施，强化消毒措施。中小学校和托幼机构切实落实新生入学查验预防接种证制度、晨检午检制度、因病缺勤病因追踪与登记制度。
    \item \imp{完善应急预案体系}：应包含总则、组织机构与职责、监测预警、信息报告、应急响应、后期处置、保障措施、应急培训和演练等基本要素，并对预案实行动态管理。
    \item \imp{改善学校卫生环境}：按照国家标准的规定，提供合乎安全卫生要求的教学环境和生活条件。
    \item \imp{定期开展健康教育}：《学校卫生工作条例》明确学校卫生任务之一就是进行学校健康教育，"传染病防控宣传教育制度"是重要制度之一。
    \item \imp{提升应急响应能力}：与相关卫生部门合作，主动加强预警预测工作，加大对校内疫情的监测力度，定期开展演练。
\end{enumerate}

\textbf{（二）学校传染病事件应急处置}

\textbf{1. 疫情报告：}

中小学校设立学校传染病疫情报告人（School Epidemic Information Whistleblower），即负责传染病疫情报告的学校专职或者兼职卫生专业技术人员、保健教师，或经培训合格的学校其他在编人员。

当发生法定传染病疫情或突发公共卫生事件时，学校疫情报告人应在24小时内向属地疾病预防控制机构和教育行政部门报告相关信息，鉴别依据为：在同一宿舍或者同一班级，1天内有3例或者连续3天内有多个学生（5例以上）患病，并有相似症状（如发热、皮疹、腹泻、呕吐、黄疸等）或者共同用餐、饮水史；个别学生出现不明原因的高热、呼吸急促或剧烈呕吐、腹泻等症状；学校发生群体性不明原因疾病或者其他突发公共卫生事件。

\textbf{2. 疫情调查：}学校全面配合卫生部门开展疫情分析、病例诊治以及流行病学调查和疫情处理工作。

\textbf{3. 控制疫情：}
\begin{itemize}[leftmargin=*]
    \item 及时隔离，安排就诊：对确诊患有法定传染病的学生、疑似病人或传染病密切接触者，配合卫生部门依法进行隔离或者医学观察，并及时安排就诊，做好检疫期相关记录。
    \item 疫情防控合作：配合属地疾病预防控制机构开展消毒、疫情调查和宣传教育等工作。
    \item 应急疫苗接种与预防用药：根据需要组织开展应急疫苗接种、预防性服药。
    \item 复课：学生病愈且隔离期满时，应到学校医务室或卫生室查验后方可进班复课。
    \item 停课并暂停集会：在传染病暴发、流行时，学校应停止举办大型师生集会和会议，采取临时停课或暂时关闭措施。
    \item 做好个人防护：教职员工在照顾患病学生、接触可能受到污染的物品或排泄物时，应根据实际情况采取必要的个人防护措施。
\end{itemize}

\textbf{4. 信息发布：}教育部门和学校要按照有关规定作好信息发布工作，信息发布要及时主动、实事求是，准确把握、注意技巧，正确引导舆论，维护社会稳定。

\textbf{5. 健康教育：}学校应根据事件性质，有针对性地开展相关传染病预防知识的教育活动，提高师生健康意识和自我防护能力，开展心理危机干预工作以消除学生的情绪波动和心理障碍。

\textbf{6. 应急反应的终止：}学校传染病事件应急反应终止的条件是末例传染病病例发生后经过最长潜伏期无新的病例出现。由卫生行政部门对学校传染病事件作出应急反应的终止决定后，学校方可解除应急措施。

\textbf{7. 评估总结：}学校在传染病暴发、流行事件得到控制后，对事件进行评估总结，评估内容主要包括事件概况、现场调查处理概况、病人救治情况、所采取措施的效果评价、应急处理过程中存在的问题和取得的经验及改进建议。评估报告须上报教育行政部门。
\end{keybox}

% ----- 12.3 食源性疾病与食品安全 -----
\section{学校食源性疾病与食品安全监督}

\subsection{学校食源性疾病和食品安全事故}

\begin{defbox}
\textbf{食源性疾病（Foodborne Disease）：}根据《中华人民共和国食品安全法》（2021年修订），食源性疾病是指食品中致病因素进入人体引起的感染性、中毒性等疾病，包括食物中毒。

\textbf{食物中毒（Food Poisoning）：}是指食用了被有毒有害物质污染的食品或者食用了含有毒有害物质的食品后出现的急性、亚急性疾病，不包括因暴饮暴食而引起的急性胃肠炎、食源性肠道传染病和寄生虫病，也不包括因一次大量或者长期少量多次摄入某些有毒、有害物质而引起的以慢性毒害为主要特征（如致畸、致癌、致突变）的疾病。

\textbf{食品安全事故（Food Safety Incidents）：}是指食源性疾病、食品污染等源于食品，对人体健康有危害或者可能有危害的事故。学校要加强食品安全管理，配合卫生健康部门加强学校食品安全监督。
\end{defbox}

\subsection{学校食物中毒事件的类型和流行病学特点}

\begin{keybox}
\textbf{学校食物中毒事故的类型：}
\begin{itemize}[leftmargin=*]
    \item \imp{重大学校食物中毒事故}：一次中毒100人以上并出现死亡病例，或出现10例及以上死亡病例
    \item \imp{较大学校食物中毒事故}：一次中毒100人及以上，或出现死亡病例
\end{itemize}

\textbf{学校食物中毒的病因分型：}细菌性食物中毒、真菌性食物中毒、有毒动植物中毒、化学性食物中毒。

\textbf{学校食物中毒事故的流行特征：}呈现季节性波动；存在地区和学校分布差异；致病因素主要为微生物和化学物质。
\end{keybox}

\subsection{学校食品安全与监督}

\begin{defbox}
\textbf{学校的职责：}食堂建筑、设备与环境；食品采购与储存；食品加工；食堂从业人员；管理与监督。

\textbf{教育部门的职责：}加强行政管理；组织人员培训；定期检查督促。

\textbf{卫生管理部门的职责：}加强卫生监督；严格卫生许可工作的审核管理和督查力度。
\end{defbox}

\subsection{学校食品安全事故的处置}

\begin{keybox}
\begin{enumerate}[leftmargin=*]
    \item \imp{报告}：2小时内报告当地教育行政部门和卫生行政部门；报告的主要内容包括发生食物中毒暴发事件单位、地点、时间、中毒人数、主要临床症状等。
    \item \imp{封存、销毁}：对疑似导致食品安全事故的食品及其原料立即封存，并由卫生机构进行检验；对确认属于被污染的食品及其原料，立即停止销售，予以召回和销毁。封存被污染的食品用工具及用具，并进行彻底清洗消毒。
    \item \imp{迅速组织救治}：发生学校食物中毒事故后，教育行政部门和学校要把治病救人放在首位，迅速组织救治，尤其是危重患者，要不惜一切代价，全力抢救。
    \item \imp{配合调查和善后}：学校应配合卫生行政部门进行调查，按卫生行政部门的要求如实提供有关材料和样品。
    \item \imp{尽早恢复正常的教学秩序}：学校食物中毒事故应急处置完成后，学校应尽快恢复正常的教学秩序。
    \item \imp{正确发布相关信息}：信息发布要及时主动、准确把握，实事求是，正确引导舆论，注重社会效果。
\end{enumerate}
\end{keybox}

% ----- 12.4 群体心因性反应 -----
\section{学校群体心因性反应事件的应对}

\subsection{群体心因性疾病概述}

\begin{keybox}
\textbf{群体心因性疾病（Mass Psychogenic Illness）：}亦称流行性癔症（Epidemic Hysteria）、群体性癔症（Mass Hysteria）或群体社会性疾病（Mass Sociogenic Illness），是指在一个紧密联系的群体中，由心理因素（如恐惧、焦虑、暗示等）引发的多人转换障碍（Conversion Disorder），该种精神障碍常表现为具有躯体症状但没有明确的器质性病变。根据症状特征，可分为焦虑型癔症（Anxiety Hysteria）和运动型癔症（Motor Hysteria）。

\textbf{触发因素：}疾病暴发通常是由感官体验引发的，如难闻的气味、可疑的物体或物质、不明的声音或其他形式的刺激，这种刺激叫做触发因素或触发因子（Trigger）。学校群体心因性疾病暴发的触发因素通常是以公共卫生事件为主，其中尤以集体预防接种和服药、疑似食物中毒最为常见，而迷信谣言、心理压力和情绪紧张等因素也常有报道。
\end{keybox}

\subsection{流行特征}

\begin{keybox}
\textbf{（一）首发病例（指示病例）：}发病与紧张或恐惧感有关，少数有"偶合症"，如偶合上呼吸道感染或其他器质性病变。

\textbf{（二）临床特征：}常见的为焦虑型癔症，常见症状包括腹痛、头痛、头晕、晕厥、恶心和换气过度；运动型癔症少见，常见症状包括舞蹈、抽搐、颤抖、行走困难、无法控制的大笑和哭泣、沟通困难和恍惚状态等。未发现有器质性病变。

\textbf{（三）三间分布：}
\begin{itemize}[leftmargin=*]
    \item \imp{地区分布}：全球范围内均有报道，多发生于经济文化落后的地区，农村地区罹患率高于城市。
    \item \imp{人群分布}：儿童青少年普遍易感；首发病例和续发病例多为女生；小学生占比较高，罹患率高于中学、职高和大学；以怀疑食物中毒和水体污染为触发因素的事件罹患率较高。
    \item \imp{时间分布}：全年均有可能暴发，但以春末夏初、夏末秋初最常见；首发病例潜伏期多在30分钟以内，续发病例潜伏期大多在首发病例发生后1天内。
\end{itemize}

\textbf{（四）多变性：}大规模群体心因性疾病往往是社会问题的突出反映，其会随着时代背景的变化而呈现多变性特点。20世纪之前主要以运动型癔症为主；20世纪以来主要以焦虑型癔症为主；从20世纪80年代初到21世纪初，化学和生物恐怖主义突起；2022年德国报道了由一名图雷特综合征网红引发的群体心因性反应的大范围暴发，证实了社交媒体作为群体心因性暴发的新媒介。
\end{keybox}

\subsection{早期识别与应急处置}

\begin{keybox}
\textbf{（一）暴发调查}

群体心因性疾病的暴发调查需要按照突发事件十步调查法开展调查：通过描述性流行病学研究查明流行分布形成病因假设；绘制发病曲线判断是同源性一次暴发还是非同源性数次暴发；确定人群分布（性别比例、年龄范围等）；确定场所分布（共同暴露因素及暴露与发病的相关性）；检验病因假设；结合实验室结果排除其他诊断，明确群体心因性疾病诊断。

\textbf{（二）识别诊断——陶芳标等构建的指标体系：}

\begin{enumerate}[leftmargin=*]
    \item \imp{诱因}：无明确病因（如未食用某种食物或药物、未吸入某种气体、未预防接种等）；或有报告有引发的间接因素如环境拥挤、通气不良、潮湿等导致疲劳、虚弱。
    \item \imp{发病者特征}：首发病例和续发病例多见于女性；学习成绩不良、工作表现不佳、经常受人批评、经常受人欺压的儿童少年易感性相对较高；病例最初出现在有相同生活或文化背景的人群中，未接触某种疑似病因者也有发病。
    \item \imp{事件发展过程}：有通过他人暗示而相继发病者；发生的时间相对集中，症状出现和消失时间快（一发都发，一停都停）；停止接触某种疑似病因后新病例仍继续出现；发病以"离心"的趋势向外发展；病例非单峰分布。
    \item \imp{临床诊断治疗}：主诉、临床症状和体征与实验室或体格检查结果不相符；非特异性治疗（如隔离、精神安慰、心理疏导或注射维生素C等非特异性药物）后，症状很快好转或消失。
\end{enumerate}

\textbf{（三）现场应急处置}

\begin{enumerate}[leftmargin=*]
    \item \imp{隔离患者}：立即将首发、继发病例转移出现场，适当隔离，以避免患者间互相影响及效仿，减少症状的顽固性和丰富性。
    \item \imp{消除紧张性情绪环境}：(1)在隔离患者的同时，引导其他学生也从该环境撤离；(2)消除来自周围环境的不良的言语暗示、动作暗示；(3)医生认真做详细检查，解释病情，使学校领导、教师、家属和围观者对治疗建立信心；(4)待症状缓解后，帮助患者分析发病的主客观原因，指导他们和家属解除相关不良精神因素。
    \item \imp{对症治疗}：对患者表现出的躯体症状采用止吐、止泻、止痛等对症治疗；对少数精神反应特别大的个体可适量使用镇静药物。
\end{enumerate}

\textbf{（四）心理治疗}

\begin{enumerate}[leftmargin=*]
    \item \imp{移情解释法}：采用符合患者心理要求的内容及形式，鼓励其参加感兴趣的活动以转移注意力。
    \item \imp{暗示疗法}：有语言暗示、药物暗示、催眠疗法等。
    \item \imp{着重治疗关键患者（首发患者）}：通过重点治疗使其痊愈，再通过其现身说法，对其他患者产生正向影响，对控制事件发展起事半功倍的作用。
    \item \imp{争取家长配合}：学校群体心因性疾病的触发因素并非直接影响儿童，而是经过家长等成人的过滤、整合后传递给儿童成为暗示信号。在预防、应急处理和后续的心理干预阶段均应重视对家长的教育引导解释。
\end{enumerate}
\end{keybox}

% ----- 复习题 -----
\section{复习题}

\begin{enumerate}[leftmargin=*, itemsep=8pt]
    \item \textbf{【名词解释】}学校突发公共卫生事件；学校传染病事件；食源性疾病；食物中毒
    \item \textbf{【名词解释】}群体心因性疾病；触发因素（触发因子）；学校传染病疫情报告人
    \item \textbf{【简答题】}简述学校突发公共卫生事件的类型和分级。
    \item \textbf{【简答题】}简述学校传染病事件的流行病学特点。
    \item \textbf{【简答题】}简述学校传染病监督管理的主要内容。
    \item \textbf{【简答题】}简述学校传染病事件应急处置的流程（疫情报告→调查→控制→终止→评估）。
    \item \textbf{【简答题】}简述学校食品安全事故的处置要点。
    \item \textbf{【简答题】}简述群体心因性疾病的定义、临床特征和三间分布。
    \item \textbf{【简答题】}简述群体心因性疾病的识别诊断指标体系（陶芳标等构建）。
    \item \textbf{【简答题】}简述群体心因性疾病的现场应急处置和心理治疗方法。
\end{enumerate}

\subsection*{参考答案要点}

\begin{enumerate}[leftmargin=*, itemsep=5pt]
    \item \textbf{学校突发公共卫生事件}：在学校内突然发生，造成或可能造成师生员工健康严重损害的重大传染病疫情、重大食物或职业中毒、不明原因的群体性疾病、群体性异常反应等事件。\textbf{学校传染病事件}：学校某种传染病在短时间内发生、波及范围广泛，发病率远超常年水平的紧急情况，占学校突发公共卫生事件的90\%$\sim$95\%。\textbf{食源性疾病}：食品中致病因素进入人体引起的感染性、中毒性等疾病，包括食物中毒。\textbf{食物中毒}：食用了被有毒有害物质污染的食品或含有毒有害物质的食品后出现的急性、亚急性疾病。

    \item \textbf{群体心因性疾病}：在一个紧密联系的群体中，由心理因素（恐惧、焦虑、暗示等）引发的多人转换障碍，表现为具有躯体症状但没有明确器质性病变。\textbf{触发因素}：引发群体心因性疾病暴发的感官体验（难闻气味、可疑物体、不明声音等），以集体预防接种和服药、疑似食物中毒最为常见。\textbf{学校传染病疫情报告人}：负责传染病疫情报告的学校专职或兼职卫生专业技术人员、保健教师，或经培训合格的学校其他在编人员。

    \item 类型：重大传染病疫情、预防接种和预防服药群体性不良事件、群体性不明原因疾病、食物中毒、其他中毒、环境因素事件、意外辐射照射事件。分级（2006年预案）：I级（特别重大）、II级（重大）、III级（较大）、IV级（一般）。

    \item (1)多以呼吸道和消化道传染病为主（两者占总事件的97.91\%）；(2)存在地区和学校分布差异（西南省份多发，乡村中小学和县城小学最多）；(3)全年两个发病高峰（3$\sim$6月和10$\sim$12月之间双峰分布）。

    \item 六个方面：(1)建设卫生管理制度（制定完善制度+实施方案+监督机制+经费人员+专业培训）；(2)标准化卫生工作流程（饮用水食堂管理+洗手消毒+查验接种证+晨检午检+因病缺勤追踪）；(3)完善应急预案体系（总则+机构职责+监测预警+信息报告+应急响应+后期处置+保障措施+培训演练+动态管理）；(4)改善学校卫生环境；(5)定期开展健康教育；(6)提升应急响应能力。

    \item 七步流程：(1)疫情报告（24小时内向疾控机构和教育行政部门报告，依据：1天3例或连续3天5例以上相似症状等）；(2)疫情调查（配合卫生部门）；(3)控制疫情（隔离就诊+防控合作+应急接种/预防用药+复课查验+停课暂停集会+个人防护）；(4)信息发布（及时主动、实事求是、正确引导舆论）；(5)健康教育（针对性防病知识+心理危机干预）；(6)应急反应终止（末例病例后经过最长潜伏期无新病例）；(7)评估总结（事件概况+调查处理+救治+措施效果+问题经验+上报教育行政部门）。

    \item 六项要点：(1)报告（2小时内报告当地教育行政部门和卫生行政部门）；(2)封存销毁（封存可疑食品及原料，确认污染者召回销毁，封存被污染工具并清洗消毒）；(3)迅速组织救治（治病救人放首位，危重患者全力抢救）；(4)配合调查和善后（如实提供材料和样品）；(5)尽早恢复正常的教学秩序；(6)正确发布相关信息（及时主动、准确把握、实事求是、正确引导舆论）。

    \item 定义见第2题。临床特征：焦虑型癔症（腹痛、头痛、头晕、晕厥、恶心、换气过度）常见，运动型癔症（舞蹈、抽搐、颤抖、行走困难、无法控制大笑哭泣、沟通困难、恍惚状态）少见，未发现器质性病变。三间分布：地区——全球均有，经济文化落后地区多发，农村$>$城市；人群——儿童青少年普遍易感，首/续发病例多为女生，小学生占比高；时间——全年均可能，春末夏初/夏末秋初最常见，首发潜伏期多在30分钟内，续发大多在首发后1天内。

    \item 陶芳标等构建的四方面指标体系：(1)诱因——无明确病因（未食用可疑食物/药物、未吸入某种气体等），或有间接因素（环境拥挤、通气不良、潮湿等导致疲劳虚弱）；(2)发病者特征——多见于女性、学习成绩不良/经常受批评欺压者易感性高、病例最初出现在相同生活/文化背景人群中、未接触疑似病因者也发病；(3)事件发展过程——通过他人暗示相继发病、时间集中症状出现和消失快（一发都发一停都停）、停止接触疑似病因后新病例仍出现、发病以"离心"趋势向外发展、病例非单峰分布；(4)临床诊断治疗——主诉/症状/体征与实验室/体格检查结果不相符、非特异性治疗后症状很快好转或消失。

    \item 现场应急处置：(1)隔离患者（转移出现场适当隔离，避免互相影响效仿）；(2)消除紧张性情绪环境（引导其他学生撤离+消除不良言语/动作暗示+医生详细检查解释建立信心+症状缓解后帮助分析主客观原因）；(3)对症治疗（止吐止泻止痛等+必要时少量镇静药物）。心理治疗：(1)移情解释法（鼓励参加感兴趣活动转移注意力）；(2)暗示疗法（语言暗示、药物暗示、催眠疗法等）；(3)着重治疗关键患者/首发患者（通过现身说法产生正向影响，事半功倍）；(4)争取家长配合（触发因素经家长过滤整合后传递给孩子成为暗示信号，需重视家长教育引导）。
\end{enumerate}

\newpage
